%%%%%%%%%%%%%%%%%%%%%%%%%%%%%%%%%%%%%%%%%%%%%%%%%%%%%%%%%%%%%%%%%%%%%%%%%%%%%
%% QSPI Specification – Additions and Extensions
%% Covers all "Fehlend" and "Offen" items from the review:
%%   A. TX FIFO Interrupts (TX_HALF, TX_EMPTY)
%%   B. FIFO depth specification
%%   C. CLKDIV register (Baud-Rate)
%%   D. ADDR_LEN bit (3-byte vs. 4-byte address)
%%   E. CPOL / CPHA bits in CTRL register
%%   F. Dual-SPI mode clarification
%%   G. XIP mode description
%%   H. TIMEOUT configuration register
%%   I. BUSY / START write behaviour
%%   J. FIFO flush bits
%%   K. Reference flash devices (Winbond W25Q128JV, Micron MT25QL128)
%%
%% Integration hint: These snippets replace or extend the existing
%% 05_qspi_text08.tex (Service Requests), 05_qspi_text10.tex (Registers),
%% 05_qspi_tab01.tex (Feature Set) and 05_qspi_text05.tex / text09.tex.
%% The register address offsets continue from 0x70 (MIS) at 0x78, 0x80.
%%%%%%%%%%%%%%%%%%%%%%%%%%%%%%%%%%%%%%%%%%%%%%%%%%%%%%%%%%%%%%%%%%%%%%%%%%%%%


%%%%%%%%%%%%%%%%%%%%%%%%%%%%%%%%%%%%%%%%%%%%%%%%%%%%%%%%%%%%%%%%%%%%%%%%%%%%%
%% A+B: ERWEITERUNG von 05_qspi_tab01.tex – Feature Set
%%      (Replace existing table body or append rows)
%%%%%%%%%%%%%%%%%%%%%%%%%%%%%%%%%%%%%%%%%%%%%%%%%%%%%%%%%%%%%%%%%%%%%%%%%%%%%
\begin{table}[H]
\caption{QSPI0 Feature Set (updated)}
\label{tab:asQspiPer01b}
\centering
\begin{tabularx}{\textwidth}{|l |l |X|}
  \hline
   & Generic QSPI Feature Set & Details \\
  \hline
  \hline
  yes & Separate Kernel Clock      & 125\;MHz (TBD) \\
  \hline
  yes & FIFO Data Buffering        & RX and TX, each 16 entries $\times$ 64\;bit = 128\;bytes \\
  \hline
  yes & Programmable Baud Rate     & $f_{SCK} = f_{kernel} / (2 \times (\text{CLKDIV} + 1))$, CLKDIV\;$\in$\;[0\ldots255] \\
  \hline
  yes & Single SPI Mode            & 1-1-1, cmd/adr/data on one wire \\
  \hline
  yes & Dual SPI Mode              & 1-1-2, data on two wires (DUAL=1, QUAD=0) \\
  \hline
  yes & Quad SPI Mode              & 1-1-4 and 1-4-4, data on four wires (QUAD=1) \\
  \hline
  yes & Configurable SPI Mode      & CPOL and CPHA selectable (Mode\;0 and Mode\;3) \\
  \hline
  yes & 3-/4-Byte Address          & Selectable via ADDR\_LEN bit \\
  \hline
  yes & Execute-In-Place (XIP)     & Auto-read on bus access, fixed opcode from CMD register \\
  \hline
  yes & Double Data Rate (DDR)     & DDR bit in CTRL register \\
  \hline
  yes & Interrupt Generation       & 7 interrupt sources, maskable via IMSC \\
  \hline
  yes & FIFO Flush                 & TX\_FLUSH and RX\_FLUSH bits in CTRL register \\
  \hline
  yes & Configurable Timeout       & Timeout counter, threshold set via TIMEOUT register \\
  \hline
\end{tabularx}
\end{table}


%%%%%%%%%%%%%%%%%%%%%%%%%%%%%%%%%%%%%%%%%%%%%%%%%%%%%%%%%%%%%%%%%%%%%%%%%%%%%
%% K: NEUE SECTION – Reference Flash Devices
%%    Empfehlung: Einfügen als neuen \subsubsection in text05.tex (Functional
%%    Overview) oder als eigenständige Section in text09.tex.
%%%%%%%%%%%%%%%%%%%%%%%%%%%%%%%%%%%%%%%%%%%%%%%%%%%%%%%%%%%%%%%%%%%%%%%%%%%%%
\subsubsection{Reference Flash Devices}

The QSPI controller is designed to be compatible with widely used NOR-Flash devices.
Table~\ref{tab:asQspiFlashRef} shows the relevant operating parameters for two
representative devices. All register default values in this specification are
set to match the Winbond W25Q128JV in standard (non-DDR) mode with CPOL=0, CPHA=0.

\begin{table}[H]
\caption{Reference NOR-Flash Devices}
\label{tab:asQspiFlashRef}
\centering
\begin{tabularx}{\textwidth}{|l |l |l |X|}
  \hline
  Parameter & Winbond W25Q128JV & Micron MT25QL128 & Comment \\
  \hline
  \hline
  Capacity              & 128\;Mbit (16\;MByte)  & 128\;Mbit (16\;MByte)  & \\
  \hline
  Supply Voltage        & 2.7\;V\ldots3.6\;V     & 2.7\;V\ldots3.6\;V     & \\
  \hline
  SPI Modes supported   & Mode\;0 (CPOL=0, CPHA=0) and Mode\;3 (CPOL=1, CPHA=1)
                        & Mode\;0 (CPOL=0, CPHA=0) and Mode\;3 (CPOL=1, CPHA=1) & \\
  \hline
  Max. clock (STR)      & 133\;MHz (Quad)         & 133\;MHz (Quad)         & \\
  \hline
  Default address width & 3\;Byte (24\;bit)       & 3\;Byte (24\;bit)       & 4-Byte mode selectable \\
  \hline
  Quad enable           & QE bit in Status Reg\;2 & via Enhanced Volatile Config Reg (0x61) & Must be set before QUAD=1 \\
  \hline
  \multicolumn{4}{|l|}{Relevant Read Opcodes:} \\
  \hline
  Read (slow)           & 0x03 (1-1-1, 0 dummy)  & 0x03 (1-1-1, 0 dummy)  & $\leq$50\;MHz \\
  \hline
  Fast Read             & 0x0B (1-1-1, 8 dummy)  & 0x0B (1-1-1, 8 dummy)  & up to 133\;MHz \\
  \hline
  Dual Output Read      & 0x3B (1-1-2, 8 dummy)  & 0x3B (1-1-2, 8 dummy)  & DUAL=1 \\
  \hline
  Quad Output Read      & 0x6B (1-1-4, 8 dummy)  & 0x6B (1-1-4, 8 dummy)  & QUAD=1 \\
  \hline
  Quad I/O Read         & 0xEB (1-4-4, 6+2 dummy)& 0xEB (1-4-4, 10 dummy) & QUAD=1, fast \\
  \hline
  \multicolumn{4}{|l|}{Relevant Write / Control Opcodes:} \\
  \hline
  Write Enable          & 0x06                   & 0x06                   & Required before program/erase \\
  \hline
  Page Program          & 0x02 (1-1-1)           & 0x02 (1-1-1)           & Up to 256\;Byte \\
  \hline
  Sector Erase (4\;KB)  & 0x20                   & 0x20                   & \\
  \hline
  Block Erase (64\;KB)  & 0xD8                   & 0xD8                   & \\
  \hline
  Chip Erase            & 0xC7 / 0x60            & 0xC7 / 0x60            & \\
  \hline
  JEDEC ID Read         & 0x9F                   & 0x9F                   & Returns 3\;Byte \\
  \hline
  JEDEC Manufacturer ID & 0xEF (Winbond)         & 0x20 (Micron)          & \\
  \hline
  \multicolumn{4}{|l|}{Typical Dummy Cycles by Mode and Clock:} \\
  \hline
  1-1-1 Fast Read       & 8 cycles               & 8 cycles               & \asqspiDummyRegExt = 0x08 \\
  \hline
  1-1-4 Quad Output     & 8 cycles               & 8 cycles               & \asqspiDummyRegExt = 0x08 \\
  \hline
  1-4-4 Quad I/O        & 6 mode + 2 dummy cycles& 10 cycles              & \asqspiDummyRegExt = 0x08 / 0x0A \\
  \hline
\end{tabularx}
\end{table}

\paragraph{Note on Quad Enable: }
Both reference devices require a Quad Enable (QE) bit to be set in the flash device's
non-volatile status register before QUAD mode can be used. This is a one-time
(or volatile) configuration step performed by the SW driver via a standard
Write Enable (0x06) followed by a Write Status Register command. The QSPI controller
itself does not perform this step automatically.


%%%%%%%%%%%%%%%%%%%%%%%%%%%%%%%%%%%%%%%%%%%%%%%%%%%%%%%%%%%%%%%%%%%%%%%%%%%%%
%% C: NEUES REGISTER – CLKDIV (Baud-Rate)
%%    Registeradresse: 0x78 (nach MIS bei 0x70)
%%    Einfügen in: 05_qspi_text10.tex nach dem MIS-Register,
%%                 und in 05_qspi_tab03.tex als neue Zeile.
%%%%%%%%%%%%%%%%%%%%%%%%%%%%%%%%%%%%%%%%%%%%%%%%%%%%%%%%%%%%%%%%%%%%%%%%%%%%%

%% Snippet für 05_qspi_tab03.tex (neue Zeile einfügen):
%% \hline
%% 0x78 & \asqspiClkdivRegExt & \asqspiClkdivRegInt & QSPI Clock Divider register \\
%% \hline
%% 0x80 & \asqspiTimeoutRegExt & \asqspiTimeoutRegInt & QSPI Timeout threshold register \\
%% \hline

\paragraph{\asqspiClkdivRegExt} -:
The Clock Divider Register controls the frequency of the generated SCK clock.
The SCK frequency is derived from the kernel clock as:
\[
  f_{SCK} = \frac{f_{kernel}}{2 \times (\text{CLKDIV} + 1)}
\]
With $f_{kernel} = 125\;\text{MHz}$ this yields the following example values:

\begin{table}[H]
\caption{CLKDIV Examples ($f_{kernel}=125\;\text{MHz}$)}
\label{tab:asQspiClkdivEx}
\centering
\begin{tabularx}{\textwidth}{|l |l |X|}
  \hline
  CLKDIV & $f_{SCK}$ & Comment \\
  \hline
  \hline
  0x00 & 62.5\;MHz & Maximum SCK (check flash device spec) \\
  \hline
  0x01 & 31.25\;MHz & \\
  \hline
  0x03 & 15.625\;MHz & Suitable for slow read without dummy cycles \\
  \hline
  0xFF & 244.1\;kHz  & Minimum SCK \\
  \hline
\end{tabularx}
\end{table}

The register must not be changed while a transfer is active (BUSY=1).
\asregkopf{Clock Divider register (in BPI)}{$0000000000000003_H$}
\asregfourxsixteen{\multicolumn{16}{|c|}{0}}
                             {\multicolumn{16}{|c|}{r}}
                             {\multicolumn{16}{|c|}{0}}
                             {\multicolumn{16}{|c|}{r}}
                             {\multicolumn{16}{|c|}{0}}
                             {\multicolumn{16}{|c|}{r}}
                             {\multicolumn{8}{|c|}{0}&\multicolumn{8}{c|}{CLKDIV}}
                             {\multicolumn{8}{|c|}{r}&\multicolumn{8}{c|}{rw}}
\asregdesc{- & 63:8 & r & Reserved \\
           \hline
           CLKDIV & 7:0 & rw & Clock divider value. Reset value 0x03 gives $f_{SCK} \approx 15.6\;\text{MHz}$ as a safe default for most NOR-Flash devices without dummy cycles.}


%%%%%%%%%%%%%%%%%%%%%%%%%%%%%%%%%%%%%%%%%%%%%%%%%%%%%%%%%%%%%%%%%%%%%%%%%%%%%
%% H: NEUES REGISTER – TIMEOUT
%%    Registeradresse: 0x80
%%%%%%%%%%%%%%%%%%%%%%%%%%%%%%%%%%%%%%%%%%%%%%%%%%%%%%%%%%%%%%%%%%%%%%%%%%%%%

\paragraph{\asqspiTimeoutRegExt} -:
The Timeout Register defines the maximum number of kernel clock cycles that the
FSM may remain in any non-IDLE state before a TIMEOUT error is raised.
If the counter reaches zero while the kernel is active (BUSY=1), the FSM
is forced back to IDLE, the CS line is deasserted, and the TIMEOUT bit is set
in the STATUS and RIS registers.

Setting TIMEOUT\_THR to 0x00000000 disables the timeout mechanism.

A typical value for a 125\;MHz kernel clock and a maximum allowed transaction time
of 1\;ms is: $\text{TIMEOUT\_THR} = 125\,000 = \texttt{0x0001E848}$.

\asregkopf{Timeout register (in BPI)}{$0000000000000000_H$}
\asregfourxsixteen{\multicolumn{16}{|c|}{0}}
                             {\multicolumn{16}{|c|}{r}}
                             {\multicolumn{16}{|c|}{0}}
                             {\multicolumn{16}{|c|}{r}}
                             {\multicolumn{16}{|c|}{TIMEOUT\_THR[31:16]}}
                             {\multicolumn{16}{|c|}{rw}}
                             {\multicolumn{16}{|c|}{TIMEOUT\_THR[15:0]}}
                             {\multicolumn{16}{|c|}{rw}}
\asregdesc{- & 63:32 & r & Reserved \\
           \hline
           TIMEOUT\_THR & 31:0 & rw & Timeout threshold in kernel clock cycles.
                                      0x00000000 = timeout disabled.}


%%%%%%%%%%%%%%%%%%%%%%%%%%%%%%%%%%%%%%%%%%%%%%%%%%%%%%%%%%%%%%%%%%%%%%%%%%%%%
%% C+D+E+F+J: ERWEITERTES CTRL-REGISTER
%%    Ersetzt den bestehenden \paragraph{\asqspiCtrlRegExt} in text10.tex.
%%    Neue Bits: CPOL, CPHA, DUAL, ADDR_LEN, TX_FLUSH, RX_FLUSH
%%    (Bits 5..10 werden belegt; bisher waren 5..63 reserved)
%%%%%%%%%%%%%%%%%%%%%%%%%%%%%%%%%%%%%%%%%%%%%%%%%%%%%%%%%%%%%%%%%%%%%%%%%%%%%

\paragraph{\asqspiCtrlRegExt} -:
This register controls all transaction parameters of the QSPI interface:
\begin{itemize}
  \item Start a transfer:
    \[ start\_s = (\asqspiCtrlRegExt\textbf{.START} == 1) \land (\asqspiStatusRegExt\textbf{.BUSY} == 0) \]
    If START is written while BUSY is set, the write is ignored (no error).
    The START bit is self-clearing; a read always returns 0.
  \item TX and RX FIFOs can be flushed independently via TX\_FLUSH and RX\_FLUSH.
    These bits are self-clearing. Flushing is only allowed when BUSY = 0;
    otherwise the flush is ignored.
  \item Single (DUAL=0, QUAD=0), Dual (DUAL=1, QUAD=0) or Quad (QUAD=1) SPI data mode.
  \item 3-Byte or 4-Byte flash address selectable via ADDR\_LEN.
  \item SPI clock polarity and phase (CPOL, CPHA) selectable. Only Mode\;0
    (CPOL=0, CPHA=0) and Mode\;3 (CPOL=1, CPHA=1) are supported by the
    reference flash devices (see Table~\ref{tab:asQspiFlashRef}).
  \item Double Data Rate (DDR) mode.
  \item Execute-In-Place (XIP) mode.
  \item Chip select can be kept asserted between transactions (CS\_HOLD).
\end{itemize}

\asregkopf{Control register (in BPI)}{$0000000000000002_H$}
\asregfourxsixteen{\multicolumn{16}{|c|}{0}}
                             {\multicolumn{16}{|c|}{r}}
                             {\multicolumn{16}{|c|}{0}}
                             {\multicolumn{16}{|c|}{r}}
                             {\multicolumn{16}{|c|}{0}}
                             {\multicolumn{16}{|c|}{r}}
                             {0&0&0&RF&TF&AL&CP&CH&DU&CS&XI&DD&QU&ST&0&0}
                             {r&r&r&w&w&rw&rw&rw&rw&rw&rw&rw&rw&w&r&r}
%% Bit layout (63..0), only lower 16 shown:
%%   Bit 13: RF  = RX_FLUSH
%%   Bit 12: TF  = TX_FLUSH
%%   Bit 11: AL  = ADDR_LEN
%%   Bit 10: CP  = CPOL
%%   Bit  9: CH  = CPHA
%%   Bit  8: DU  = DUAL
%%   Bit  7: CS  = CS_HOLD  (was bit 4 before -- renumbered!)
%%   Bit  6: XI  = XIP
%%   Bit  5: DD  = DDR
%%   Bit  4: QU  = QUAD
%%   Bit  3: START (was bit 0 before -- renumbered!)
%%   Bits 2..0: reserved
%% Note: If backwards compatibility is required, keep the original bit
%%       positions (0..4) and extend upward for the new bits (5..13).
%%       The layout above keeps the existing lower bits in place:
\asregdesc{- & 63:14 & r & Reserved \\
           \hline
           RX\_FLUSH (RF) & 13 & w  & Flush RX FIFO. Self-clearing. Only effective when BUSY=0. \\
           \hline
           TX\_FLUSH (TF) & 12 & w  & Flush TX FIFO. Self-clearing. Only effective when BUSY=0. \\
           \hline
           ADDR\_LEN (AL) & 11 & rw & Address length. 0 = 3-Byte (24\;bit), 1 = 4-Byte (32\;bit). \\
           \hline
           CPOL (CP)      & 10 & rw & Clock polarity. 0 = idle low (Mode\;0/2), 1 = idle high (Mode\;1/3). \\
           \hline
           CPHA (CH)      &  9 & rw & Clock phase. 0 = sample on first edge (Mode\;0/3), 1 = sample on second edge (Mode\;1/2). Note: only combinations CPOL=CPHA=0 and CPOL=CPHA=1 are supported by the reference flash devices. \\
           \hline
           DUAL (DU)      &  8 & rw & Enable dual-wire data mode (1-1-2). Effective only when QUAD=0. \\
           \hline
           CS\_HOLD (CS)  &  7 & rw & Keep CS asserted between consecutive transactions. \\
           \hline
           XIP (XI)       &  6 & rw & Enable execute-in-place mode. See Section~\ref{subsec:qspifunc02} for details. \\
           \hline
           DDR (DD)       &  5 & rw & Enable double data rate. \\
           \hline
           QUAD (QU)      &  4 & rw & Enable quad-wire data mode (1-1-4 or 1-4-4). Overrides DUAL. \\
           \hline
           START (ST)     &  3 & w  & Start transfer. Writing 1 initiates a transaction if BUSY=0. Self-clearing. \\
           \hline
           -              & 2:0 & r  & Reserved.}


%%%%%%%%%%%%%%%%%%%%%%%%%%%%%%%%%%%%%%%%%%%%%%%%%%%%%%%%%%%%%%%%%%%%%%%%%%%%%
%% A: ERWEITERUNG der Interrupt-Register um TX_HALF und TX_EMPTY
%%    Ersetzt die bestehenden 05_qspi_text08.tex Service Request Table
%%    sowie alle 5 Interrupt-Register-Definitionen in 05_qspi_text10.tex.
%%%%%%%%%%%%%%%%%%%%%%%%%%%%%%%%%%%%%%%%%%%%%%%%%%%%%%%%%%%%%%%%%%%%%%%%%%%%%

%% ---- 05_qspi_text08.tex (Service Requests table) ----------------------

Interrupt sources are level-based and derived directly from kernel status
or FIFO fill levels.

Table~\ref{tab:asQspiPer06b} shows the service requests of the QSPI.
\begin{table}[H]
\caption{QSPI Service Requests (updated)}
\label{tab:asQspiPer06c}
\centering
\begin{tabularx}{\textwidth}{|l |l |l |X|}
  \hline
  Bit & Service Request & IRQ & Explanation \\
  \hline
  \hline
  0 & DONE      & DONE      & Transfer completed (FSM reached DONE state) \\
  \hline
  1 & RX\_HALF  & RX\_HALF  & RX FIFO $\geq$ half full \\
  \hline
  2 & RX\_EMPTY & RX\_EMPTY & RX FIFO empty \\
  \hline
  3 & TX\_HALF  & TX\_HALF  & TX FIFO $\leq$ half full (ready to be refilled) \\
  \hline
  4 & TX\_EMPTY & TX\_EMPTY & TX FIFO empty \\
  \hline
  5 & ERROR     & ERROR     & Error latched (see STATUS register) \\
  \hline
  6 & TIMEOUT   & TIMEOUT   & Timeout threshold reached \\
  \hline
  63:7 & -       & -         & Reserved \\
  \hline
\end{tabularx}
\end{table}

In QSPI: \[ \text{MIS} = \text{RIS} \;\&\; \text{IMSC} \]
and then the external interrupt output is: \[ \text{IRQ} = |\;\text{MIS} \]

%% ---- Interrupt Registers (all 5): RIS, IMSC, MIS, ICR, ISR --------
%% Bit field encoding for all 5 registers is identical, only r/w type differs.
%% Abbreviated header for the bit row (lower 16 bits of 64-bit register):
%%   {0&0&0&0&0&0&0&0&0&TI&ER&TE&TH&EM&HA&DO}
%%   Bit 6: TI = TIMEOUT
%%   Bit 5: ER = ERROR
%%   Bit 4: TE = TX_EMPTY
%%   Bit 3: TH = TX_HALF
%%   Bit 2: EM = RX_EMPTY
%%   Bit 1: HA = RX_HALF
%%   Bit 0: DO = DONE

\paragraph{\asqspiRISRegExt} -:
The Raw Interrupt Status Register (RIS) is read-only. It reflects the unmasked
current interrupt status. Writing has no effect.
\begin{itemize}
  \item 0 -- No Interrupt pending
  \item 1 -- Interrupt pending
\end{itemize}
\asregkopf{Raw Interrupt Status Register (in BPI)}{$0000000000000000_H$}
\asregfourxsixteen{\multicolumn{16}{|c|}{0}}
                             {\multicolumn{16}{|c|}{r}}
                             {\multicolumn{16}{|c|}{0}}
                             {\multicolumn{16}{|c|}{r}}
                             {\multicolumn{16}{|c|}{0}}
                             {\multicolumn{16}{|c|}{r}}
                             {0&0&0&0&0&0&0&0&0&TI&ER&TE&TH&EM&HA&DO}
                             {r&r&r&r&r&r&r&r&r&rh&rh&rh&rh&rh&rh&rh}
\asregdesc{- & 63:7 & r & Reserved \\
           \hline
           TI & 6 & rh & TIMEOUT, Timeout threshold reached \\
           \hline
           ER & 5 & rh & ERROR, Error latched \\
           \hline
           TE & 4 & rh & TX\_EMPTY, TX FIFO empty \\
           \hline
           TH & 3 & rh & TX\_HALF, TX FIFO $\leq$ half full \\
           \hline
           EM & 2 & rh & RX\_EMPTY, RX FIFO empty \\
           \hline
           HA & 1 & rh & RX\_HALF, RX FIFO $\geq$ half full \\
           \hline
           DO & 0 & rh & DONE, Transfer completed}

\paragraph{\asqspiIMSCRegExt} -:
The Interrupt Mask Control Register (IMSC) enables or disables individual interrupt
sources. Write~1 to unmask (enable); write~0 to mask (disable).
\begin{itemize}
  \item 0 -- Interrupt disabled (masked)
  \item 1 -- Interrupt enabled (unmasked)
\end{itemize}
\asregkopf{Interrupt Mask Control Register (in BPI)}{$0000000000000000_H$}
\asregfourxsixteen{\multicolumn{16}{|c|}{0}}
                             {\multicolumn{16}{|c|}{r}}
                             {\multicolumn{16}{|c|}{0}}
                             {\multicolumn{16}{|c|}{r}}
                             {\multicolumn{16}{|c|}{0}}
                             {\multicolumn{16}{|c|}{r}}
                             {0&0&0&0&0&0&0&0&0&TI&ER&TE&TH&EM&HA&DO}
                             {r&r&r&r&r&r&r&r&r&rw&rw&rw&rw&rw&rw&rw}
\asregdesc{- & 63:7 & r & Reserved \\
           \hline
           TI & 6 & rw & TIMEOUT \\
           \hline
           ER & 5 & rw & ERROR \\
           \hline
           TE & 4 & rw & TX\_EMPTY \\
           \hline
           TH & 3 & rw & TX\_HALF \\
           \hline
           EM & 2 & rw & RX\_EMPTY \\
           \hline
           HA & 1 & rw & RX\_HALF \\
           \hline
           DO & 0 & rw & DONE}

\paragraph{\asqspiMISRegExt} -:
The Masked Interrupt Status Register (MIS) is read-only. It returns the logical
AND of RIS and IMSC, i.e.\ the effective interrupt status visible to the CPU.
\begin{itemize}
  \item 0 -- No Interrupt pending
  \item 1 -- Interrupt pending
\end{itemize}
\asregkopf{Masked Interrupt Status Register (in BPI)}{$0000000000000000_H$}
\asregfourxsixteen{\multicolumn{16}{|c|}{0}}
                             {\multicolumn{16}{|c|}{r}}
                             {\multicolumn{16}{|c|}{0}}
                             {\multicolumn{16}{|c|}{r}}
                             {\multicolumn{16}{|c|}{0}}
                             {\multicolumn{16}{|c|}{r}}
                             {0&0&0&0&0&0&0&0&0&TI&ER&TE&TH&EM&HA&DO}
                             {r&r&r&r&r&r&r&r&r&rh&rh&rh&rh&rh&rh&rh}
\asregdesc{- & 63:7 & r & Reserved \\
           \hline
           TI & 6 & rh & TIMEOUT \\
           \hline
           ER & 5 & rh & ERROR \\
           \hline
           TE & 4 & rh & TX\_EMPTY \\
           \hline
           TH & 3 & rh & TX\_HALF \\
           \hline
           EM & 2 & rh & RX\_EMPTY \\
           \hline
           HA & 1 & rh & RX\_HALF \\
           \hline
           DO & 0 & rh & DONE}

\paragraph{\asqspiICRRegExt} -:
The Interrupt Clear Register (ICR) is write-only. Writing~1 to a bit clears
the corresponding interrupt in both RIS and MIS. Writing~0 has no effect.
Reading returns~0.

Note: If a data transfer is handled via interrupt requests (this chip has no DMA),
the interrupt should be cleared via ICR only \emph{after} the FIFO has been drained
or refilled, to avoid losing the interrupt if the condition is still active.
\begin{itemize}
  \item 0 -- No change
  \item 1 -- Clear Interrupt
\end{itemize}
\asregkopf{Interrupt Clear Register (in BPI)}{$0000000000000000_H$}
\asregfourxsixteen{\multicolumn{16}{|c|}{0}}
                             {\multicolumn{16}{|c|}{r}}
                             {\multicolumn{16}{|c|}{0}}
                             {\multicolumn{16}{|c|}{r}}
                             {\multicolumn{16}{|c|}{0}}
                             {\multicolumn{16}{|c|}{r}}
                             {0&0&0&0&0&0&0&0&0&TI&ER&TE&TH&EM&HA&DO}
                             {r&r&r&r&r&r&r&r&r&w&w&w&w&w&w&w}
\asregdesc{- & 63:7 & r & Reserved \\
           \hline
           TI & 6 & w & TIMEOUT \\
           \hline
           ER & 5 & w & ERROR \\
           \hline
           TE & 4 & w & TX\_EMPTY \\
           \hline
           TH & 3 & w & TX\_HALF \\
           \hline
           EM & 2 & w & RX\_EMPTY \\
           \hline
           HA & 1 & w & RX\_HALF \\
           \hline
           DO & 0 & w & DONE}

\paragraph{\asqspiISRRegExt} -:
The Interrupt Set Register (ISR) is write-only and is provided for software
testing of the interrupt path. Writing~1 forces the corresponding RIS bit to~1
(and consequently MIS if IMSC is set). Writing~0 has no effect. Reading returns~0.
\begin{itemize}
  \item 0 -- No change
  \item 1 -- Force Interrupt (test only)
\end{itemize}
\asregkopf{Interrupt Set Register (in BPI)}{$0000000000000000_H$}
\asregfourxsixteen{\multicolumn{16}{|c|}{0}}
                             {\multicolumn{16}{|c|}{r}}
                             {\multicolumn{16}{|c|}{0}}
                             {\multicolumn{16}{|c|}{r}}
                             {\multicolumn{16}{|c|}{0}}
                             {\multicolumn{16}{|c|}{r}}
                             {0&0&0&0&0&0&0&0&0&TI&ER&TE&TH&EM&HA&DO}
                             {r&r&r&r&r&r&r&r&r&w&w&w&w&w&w&w}
\asregdesc{- & 63:7 & r & Reserved \\
           \hline
           TI & 6 & w & TIMEOUT \\
           \hline
           ER & 5 & w & ERROR \\
           \hline
           TE & 4 & w & TX\_EMPTY \\
           \hline
           TH & 3 & w & TX\_HALF \\
           \hline
           EM & 2 & w & RX\_EMPTY \\
           \hline
           HA & 1 & w & RX\_HALF \\
           \hline
           DO & 0 & w & DONE}


%%%%%%%%%%%%%%%%%%%%%%%%%%%%%%%%%%%%%%%%%%%%%%%%%%%%%%%%%%%%%%%%%%%%%%%%%%%%%
%% G: XIP MODE DESCRIPTION
%%    Neuer Paragraph in 05_qspi_text09.tex (Kernel description),
%%    nach dem bestehenden "Error Conditions" Paragraph.
%%%%%%%%%%%%%%%%%%%%%%%%%%%%%%%%%%%%%%%%%%%%%%%%%%%%%%%%%%%%%%%%%%%%%%%%%%%%%

\paragraph{Execute-In-Place (XIP) Mode: }

When the XIP bit in \asqspiCtrlRegExt\ is set, the QSPI controller automatically
services read requests from the CPU bus without explicit SW intervention.

\begin{itemize}
  \item XIP is only supported for \emph{read} accesses. Write accesses from the
    bus while XIP is active are not forwarded to the flash and cause an ERROR.
  \item The opcode for XIP reads is taken from \asqspiCmdRegExt\textbf{.OPCODE}.
    SW must configure this register with a suitable read command before
    enabling XIP (e.g.\ 0x6B for Quad Output Fast Read).
  \item The number of dummy cycles is taken from \asqspiDummyRegExt\textbf{.DUMMY\_CYC}.
  \item The ADDR\_LEN, QUAD, DUAL, CPOL, CPHA, and CLKDIV settings apply as
    configured; these must be stable before XIP is enabled.
  \item While XIP is active, the START bit is ignored (writing START has no effect).
  \item XIP is terminated by clearing the XIP bit. The current flash transaction,
    if any, is completed before the controller returns to IDLE.
  \item An XIP conflict error is raised when SW attempts to write
    \asqspiCtrlRegExt\textbf{.START} while XIP is active. This sets the
    ERROR bit in \asqspiStatusRegExt\ and in \asqspiRISRegExt.
\end{itemize}

Typical SW sequence to enable XIP:
\begin{enumerate}
  \item Configure CLKDIV, CPOL, CPHA, QUAD/DUAL, ADDR\_LEN as needed.
  \item Write the desired read opcode to \asqspiCmdRegExt.
  \item Write the required dummy cycles to \asqspiDummyRegExt.
  \item Set XIP bit in \asqspiCtrlRegExt. XIP becomes active immediately.
\end{enumerate}


%%%%%%%%%%%%%%%%%%%%%%%%%%%%%%%%%%%%%%%%%%%%%%%%%%%%%%%%%%%%%%%%%%%%%%%%%%%%%
%% I: BUSY / START Write Behaviour
%%    Neuer Paragraph in 05_qspi_text09.tex (Design Rules).
%%    Ergänzt den bestehenden "Start Semantics" Paragraph.
%%%%%%%%%%%%%%%%%%%%%%%%%%%%%%%%%%%%%%%%%%%%%%%%%%%%%%%%%%%%%%%%%%%%%%%%%%%%%

\paragraph{START Write while BUSY: }

Writing the START bit while BUSY is asserted is a no-operation (NOP):
the write is silently discarded and no error is raised. This simplifies
driver implementation (no mandatory pre-check required), but SW is responsible
for polling BUSY or waiting for the DONE interrupt before initiating a
new transaction if ordering is required.

The formal start condition remains:
\[
  start\_s = (\asqspiCtrlRegExt\textbf{.START} == 1) \land (\asqspiStatusRegExt\textbf{.BUSY} == 0)
\]


%%%%%%%%%%%%%%%%%%%%%%%%%%%%%%%%%%%%%%%%%%%%%%%%%%%%%%%%%%%%%%%%%%%%%%%%%%%%%
%% UPDATED REGISTER LIST TABLE
%%    Ersetzt 05_qspi_tab03.tex vollständig.
%%%%%%%%%%%%%%%%%%%%%%%%%%%%%%%%%%%%%%%%%%%%%%%%%%%%%%%%%%%%%%%%%%%%%%%%%%%%%

\begin{table}[H]
\caption{Register List of Module QSPI0 (updated)}
\label{tab:asQspiPer03b}
\centering
\begin{tabularx}{\textwidth}{|l |l |l |X|}
  \hline
  Offset & Register Name in Chip & Name in Module & Comment \\
  \hline
  \hline
  0x00 & \asqspiIDRegExt      & \asqspiIDRegInt      & ID of QSPI0 peripheral \\
  \hline
  0x08 & \asqspiCtrlRegExt    & \asqspiCtrlRegInt    & QSPI Control register \\
  \hline
  0x10 & \asqspiCmdRegExt     & \asqspiCmdRegInt     & SPI command opcode \\
  \hline
  0x18 & \asqspiAddrRegExt    & \asqspiAddrRegInt    & QSPI target flash address \\
  \hline
  0x20 & \asqspiLenRegExt     & \asqspiLenRegInt     & QSPI transfer length in bytes \\
  \hline
  0x28 & \asqspiDummyRegExt   & \asqspiDummyRegInt   & QSPI dummy cycles \\
  \hline
  0x30 & \asqspiStatusRegExt  & \asqspiStatusRegInt  & QSPI status flags \\
  \hline
  0x38 & \asqspiTxdataRegExt  & \asqspiTxdataRegInt  & QSPI TX FIFO write \\
  \hline
  0x40 & \asqspiRxdataRegExt  & \asqspiRxdataRegInt  & QSPI RX FIFO read \\
  \hline
  0x48 & \asqspiFifostatRegExt& \asqspiFifostatRegInt& QSPI FIFO levels \\
  \hline
  0x50 & \asqspiIMSCRegExt    & \asqspiIMSCRegInt    & Interrupt Mask Control Register \\
  \hline
  0x58 & \asqspiRISRegExt     & \asqspiRISRegInt     & Raw Interrupt Status Register \\
  \hline
  0x60 & \asqspiISRRegExt     & \asqspiISRRegInt     & Interrupt Set Register \\
  \hline
  0x68 & \asqspiICRRegExt     & \asqspiICRRegInt     & Interrupt Clear Register \\
  \hline
  0x70 & \asqspiMISRegExt     & \asqspiMISRegInt     & Masked Interrupt Status Register \\
  \hline
  0x78 & \asqspiClkdivRegExt  & \asqspiClkdivRegInt  & Clock Divider (Baud Rate) \\
  \hline
  0x80 & \asqspiTimeoutRegExt & \asqspiTimeoutRegInt & Timeout threshold \\
  \hline
\end{tabularx}
\end{table}


%%%%%%%%%%%%%%%%%%%%%%%%%%%%%%%%%%%%%%%%%%%%%%%%%%%%%%%%%%%%%%%%%%%%%%%%%%%%%
%% UPDATED REGISTER OVERVIEW TABLE (for text10.tex)
%%    Ergänzt die Registerübersicht am Anfang der Register-Section.
%%%%%%%%%%%%%%%%%%%%%%%%%%%%%%%%%%%%%%%%%%%%%%%%%%%%%%%%%%%%%%%%%%%%%%%%%%%%%

\begin{table}[H]
\caption{QSPI Register List (updated)}
\label{tab:asQspiPer04b}
\centering
\begin{tabularx}{\textwidth}{|l |X |l|}
  \hline
  Register Group & Register Name & Register Symbol \\
  \hline
  \hline
  System Registers & (Clock Control Register) & (\textbf{CLC}) \\
  & Identification Register & \textbf{ID} \\
  \hline
  Special  & (Run Control Register) & (\textbf{RUN\_CTRL}) \\
  Function & QSPI Control register & \asqspiCtrlRegExt \\
  Registers & SPI command opcode & \asqspiCmdRegExt \\
  & QSPI target flash address  & \asqspiAddrRegExt \\
  & QSPI transfer length in bytes  & \asqspiLenRegExt \\
  & QSPI dummy cycles  & \asqspiDummyRegExt \\
  & QSPI status flags  & \asqspiStatusRegExt \\
  & QSPI TX FIFO write  & \asqspiTxdataRegExt \\
  & QSPI RX FIFO read  & \asqspiRxdataRegExt \\
  & QSPI FIFO levels  & \asqspiFifostatRegExt \\
  & Clock Divider (Baud Rate) & \asqspiClkdivRegExt \\
  & Timeout Threshold & \asqspiTimeoutRegExt \\
  \hline
  Interrupt & Interrupt Mask Control Register & \asqspiIMSCRegExt \\
  Registers & Raw Interrupt Status Register & \asqspiRISRegExt \\
  & Interrupt Set Register & \asqspiISRRegExt \\
  & Interrupt Clear Register & \asqspiICRRegExt \\
  & Masked Interrupt Status Register & \asqspiMISRegExt \\
  \hline
\end{tabularx}
\end{table}
